\chapter{Discussion}
\label{ch:discussion}
This project aimed to investigate 

% The main findings of this project were:

% - Projection differences remained a good predictor of evilness and sycophancy during a sequential fine-tuning consisting of six benign supervised datasets.
% - Importantly, the regret with respect to a perfectly fitted forecaster did not increase over time for this trajectory.
% - We found evidence of the internal representations of the model drifting and not coming to normal after realignment. This relates to papers like the basin safety one. Behavior is not the same as representation drift. 
% - Despite this representation drift, a realigned Qwen2.5 experimented the same misalignment after seeing a new EM-eliciting dataset than its never-misaligned counterpart.
% - Indeed, we found that if the second EM-eliciting misalignment dataset was the same as the first one, then the model experienced less misalignment.
% - We also found a mechanistic explanation to this last findings. Updates to the model weights were smaller through training for the second training event in the same data, despite having a realignment step between those.
% - While these findings open new avenues for future research (e.g., can we create vaccines for misalignment?). We do not recommend, however, realigning a model if reinitialising the model from a checkpoint prior to misalignment is possible, because EM can be very broad and realignment could leave some misalignment leftovers as other papers in the literature suggest.

% \citet{nghiem2026trait_space} also found that benign training causes drift over long sequences. which required recalibrating the monitor. To address this problem, they added information about the training step and recalibrated the monitor using other benign long-horizon trajectories. They also reported that performance can deteriorate when the starting model changes.
% \citet{nghiem2026trait_space} also found a degradation in their monitor's performance when the starting model was already misaligned. 

In this chapter, we analyse the experiment results presented in \Cref{ch:experiments} and discuss their implications for the research questions posed in \Cref{ch:intro}. We also consider the limitations of our study and suggest directions for future work.

% mention that math datasets seem to align less the model. 
\section{RQ1: Do Projection Differences Computed From the Initial Model Remain a Good Predictor of Future Behaviour?}

\subsection{Do Projection Differences Remain Useful?}
Experiment 1 (\Cref{sec:robustness}) shows that all projection differences variants we introduced in \Cref{sec:deltap-variants}, including $\Delta P_0$, remained highly correlated with behaviour for trajectories that consisted only of normal steps. Regarding the observed differences in terms of correlation, while refreshing the persona vector's answers seemed to either hurt ot not affect the correlation, the amount of trajectories studied was small, so we cannot draw strong conclusions from this. However, we found a counterexample, meaning that a projection difference that was predictive of behaviour may stop being so. 

\subsection{RQ2: How Does Prior History Affect future misalignment?}
One of the most interesting findings of the experiments conducted was that the last dataset had the highest impact on the model's behaviour.


\section{Limitations}
\begin{itemize}
    \item We only considered one model (Qwen2.5-7B-Instruct).
    \item We only considered two traits (helpfulness and harmlessness).
    \item The length of the trajectories was relatively short (six steps). This applies to both experiments. Regarding the first experiment, it is unclear whether training on bening datasets for longer could have caused a degradation in the correlation between $\Delta P_0$ and the behaviour reached after the step.
    \item We only run one seed for each branch of Experiment 1.
    \item The correlation and prediction errors were computed from only 8 probe datasets.
    \item The main limitation of our study is that we only considered a single example dataset. More probe datasets would provide a more robust estimate of the correlation and prediction errors.
    \item 
\end{itemize}

\section{Future Work}
Future work should explor
